\section{Eliminating positive-dimensional centers}
\label{sec:centers}

The specialization comparison can be combined with a tangential
variation of a minimizing test. The essential point is that the estimates
for sections on its center are uniform as that center varies.

\begin{proposition}\label{prop:point-center}
Suppose that $P$ is not generated at $x$. There is a real number
$0<t<n+1$ such that, for arbitrarily large integers $m$, the infimum
$f_{m,t}(x)$ is attained and every minimizing monomial pair has center
exactly $\{x\}$.
\end{proposition}

Here a minimizing pair consists of a monomial test and a basis attaining
the infimum. The assertion concerns all such pairs, even if their models,
coordinates, and bases differ. This universal assertion will be needed
in Section~\ref{sec:isolation}: a second test with the same normalized
section orders and discrepancy is also minimizing, so it too must have
point center. We first establish a uniform estimate that will apply to
all possible centers.



\subsection{A uniform estimate on restriction spaces}

Fix a positive integer $k$ such that $kL$ is very ample, and use its
complete system to embed $X$ in a projective space. For an integral
$d$-dimensional closed subvariety $W\subset X$, put
\[
 I_W(h)=\dim_{\C}\im\bigl(R_h\longrightarrow H^0(W,hL|_W)\bigr).
\]
We require a lower estimate for this image, rather than surjectivity of
the restriction map.

\begin{lemma}\label{lem:uniform-restrictions}
For each $1\le d\le n$, there are constants $C_d$ and $H_d$, depending
only on $(X,L,k,d)$, such that
\begin{equation}\label{eq:uniform-restrictions}
 \bigl(d!I_W(h)\bigr)^{1/d}\ge h-C_d
 \qquad(h\ge H_d)
\end{equation}
for every integral $d$-dimensional closed subvariety $W\subset X$.
\end{lemma}

\begin{proof}
Set $D_0=k^d$. The degree of $W$ in the fixed embedding is
\[
 D=(kL)^d\cdot W=k^d(L^d\cdot W)\ge D_0,
\]
because the last intersection number is a positive integer. A general
linear projection gives a finite morphism
$\pi:W\longrightarrow\mathbb P^d$ of degree $D$. Its generic fiber is
separable, since the ground field has characteristic zero. Write
$K=\C(\mathbb P^d)$ and $E=\C(W)$.

Choose a nonzero linear coordinate $s_0$ of the projection. All points of
the geometric generic fiber lie in $s_0\ne0$. Select $D_0$ of these
distinct points. Over an algebraic closure of $K$, affine linear forms
separate any two of them. For each selected point, multiplying at most
$D_0-1$ such forms gives a polynomial vanishing at the other selected
points and not at that one. Homogenizing with $s_0$ gives ambient forms
of degree $D_0$ whose normalized values span the functions on these
$D_0$ points.

Let $V_{D_0}$ be the complex vector space of ambient degree-$D_0$
forms. Restriction and division by $s_0^{D_0}$ define a $K$-linear map
$V_{D_0}\otimes_{\C}K\longrightarrow E$. It has rank at least $D_0$.
Indeed this rank can be checked after extending scalars to the algebraic
closure of $K$, where the preceding evaluation argument applies. We may
therefore choose ambient degree-$D_0$ forms $F_1,\ldots,F_{D_0}$ for which
\[
 F_1/s_0^{D_0},\ldots,F_{D_0}/s_0^{D_0}\in E
\]
are linearly independent over $K$. The forms can be chosen among a fixed
monomial spanning set; their coefficients need not depend on elements
of $K$.

For $q\ge D_0$, multiply each $F_i$ by a monomial basis of the
homogeneous forms of degree $q-D_0$ in the coordinates of the
projection. These products restrict
to linearly independent sections on $W$. In fact, a relation between
them, divided by $s_0^q$, is a relation between the displayed
$K$-independent elements; all its coefficients must vanish in $K$.
The corresponding base forms must then vanish identically on
$\mathbb P^d$. The products are restrictions of sections in $R_{qk}$,
so
\begin{equation}\label{eq:restriction-binomial}
 I_W(qk)\ge D_0\binom{q-D_0+d}{d}.
\end{equation}
In particular,
\begin{equation}\label{eq:restriction-multiples}
 \bigl(d!I_W(qk)\bigr)^{1/d}
 \ge k(q-D_0)=qk-k^{d+1}.
\end{equation}

To cover all degrees, choose, once and for all, an integer $h_r$ in each
residue class $r$ modulo $k$ such that $h_rL$ is globally generated.
There is a section of $h_rL$ not identically zero on any given $W$.
Multiplication by that section injects the degree-$qk$ restriction space
into the degree-$(qk+h_r)$ restriction space, because $W$ is integral.
For $h=qk+h_r$ and $q\ge D_0$, Equation
\eqref{eq:restriction-multiples} thus gives
\[
 \bigl(d!I_W(h)\bigr)^{1/d}
 \ge h-(h_r+k^{d+1}).
\]
Taking maxima over the finitely many residue classes proves the lemma.
For example, both $C_d$ and $H_d$ may be taken to equal
$\max_r h_r+k^{d+1}$. All thresholds and constants are independent
of $W$.
\end{proof}

\subsection{A minimizing test and its tangent directions}

Assume from now on that $P$ is not generated at $x$.
Fix $t\in[t_0,n+1]$ and a sufficiently large degree $m$ in the ranges
of Proposition~\ref{prop:minimizers}, and write $f_j=f_{j,t}(x)$.
Suppose that a minimizing monomial pair has test $v$ and center $W$
of dimension $d>0$. We derive upper and lower bounds for its weighted
tangential moment, keeping $m$ and $t$ fixed in this calculation.
Let $Z$ be the closure of its defining stratum on
a smooth model $Y\to X$, and let $r$ be the number of its positive
weights. Put $k'=n-d$, so $r\le k'$. For the local estimates below,
keep this test, its model, and the coordinates chosen below fixed; section spaces in
every degree $j$ use the same Taylor order.

Choose a very general smooth point $y\in W$ satisfying
Lemma~\ref{lem:specialization} in degree $m$:
\begin{equation}\label{eq:center-specialization}
 f_{m,t}(y)\ge f_{m,t}(x).
\end{equation}
The very general condition and the smooth open sets used below can be
met simultaneously over $\C$. Above $y$ choose a
point $z$ of the open stratum of $Z$ where $Z\to W$ is smooth, avoiding
all additional components of the relative canonical divisor. Generic
smoothness allows these choices. Take regular parameters at $z$ in
three groups:
\begin{enumerate}
 \item the $r$ parameters defining the positive-weight crossing
 components;
 \item $k'-r$ parameters along the fiber of $Z\to W$;
 \item $d$ parameters pulled back from affine linear projective
 coordinates on $W$ vanishing at $y$.
\end{enumerate}
For the last group, choose $d$ independent differentials from the fixed
$kL$ embedding. Complete them by the fiber parameters using the
smoothness of $Z\to W$.

The point $z$ lies in the open positive-support stratum, the positive
equations change only by units, and the complementary parameters have
weight zero. By Lemma~\ref{lem:monomial-local}, the original weighted
order is therefore unchanged and can be computed from Taylor series
at $z$.

Choose a local frame of $L$ near $y$ and pull it back to $Y$. The terms
of a section with all first $r$ exponents zero are its restriction to
$Z$. Since that restriction comes from $W$, its Taylor expansion
depends only on the last $d$ parameters.

Write Taylor exponents as
\[
 (p,\beta)\in\mathbb Z_{\ge0}^{k'}\times\mathbb Z_{\ge0}^{d},
\]
and let $\ell(p)$ be their old weighted order; the middle parameters
have weight zero. Order monomials first by increasing $\ell(p)$, then
by increasing $|\beta|_1$, and finally by a fixed additive monomial
well-order, such as total degree followed by lexicographic order. This
combined order is a well-order: below any fixed weighted bound there
are only finitely many exponent choices in the positive-weight
variables; the tangential total degree and the last tie-break then
have least elements. It is also additive.

Let $\Gamma_j$ be the set of possible initial exponents of nonzero
sections of $R_j$ in these coordinates. The leading coefficient has
one-dimensional leaves, so elimination gives
\begin{equation}\label{eq:center-semigroup}
 \#\Gamma_j=N_j,
 \qquad \Gamma_j+\Gamma_q\subset\Gamma_{j+q}.
\end{equation}
The cardinality follows by the elimination argument in the proof of
Lemma~\ref{lem:min-initials}; compare the one-dimensional-leaf count in
\cite[Proposition~2.6]{kavehkhovanskii2012}.
Multiplication adds initials because the
order is additive and the local frames in all degrees are powers of
the same frame.

Take a degree-$m$ basis representing the distinct elements of
$\Gamma_m$. Since the initial order first compares $v$-values, this
basis is adapted to the $v$-filtration: its sorted valuation values
dominate those of any basis. Replacing the original minimizing basis by
this one cannot increase the objective, so the new pair still realizes
the global minimum. For this fixed test and reference degree $m$, put
\begin{equation}\label{eq:center-q-b}
 Q_j=\sum_{(p,\beta)\in\Gamma_j}e^{-\ell(p)/m}.
\end{equation}
In particular,
\begin{equation}\label{eq:center-q-minimum}
 Q_m=N_m f_m e^{-\A(v)/t}.
\end{equation}
We first use minimality to bound the weighted mean tangential degree in
$\Gamma_m$ from above.

Perturb $v$ at $z$ by giving weight $\epsilon$ to each of the $d$
tangential parameters and weight $\epsilon^2$ to each of the $k'-r$
fiber parameters, while leaving its positive weights unchanged. For
$\epsilon>0$ this is an allowed monomial test $v_\epsilon$ with center
$z$ on the model and center $y$ on $X$. Since the added coordinate
divisors have discrepancy one near $z$,
\begin{equation}\label{eq:center-discrepancy-variation}
 \A(v_\epsilon)=\A(v)+d\epsilon+(k'-r)\epsilon^2.
\end{equation}
For a section of the chosen basis with initial exponent $(p,\beta)$,
\begin{equation}\label{eq:center-order-variation}
 v_\epsilon(s)=\ell(p)+\epsilon|\beta|_1+O(\epsilon^2)
 \qquad(\epsilon\downarrow0).
\end{equation}
To justify this for Taylor series, fix the section. There are only
finitely many positive-coordinate exponent choices below any fixed
old weighted bound, so there is a positive gap from its least old
weight to the next larger weight in a bounded interval. Nonnegativity
of every exponent then excludes higher old weights from the minimum
for sufficiently small $\epsilon$. Among terms of least old weight,
the least tangential total degree determines the coefficient of
$\epsilon$; its value is $|\beta|_1$. Remaining differences affect only
the $\epsilon^2$ coefficient. One term realizing both least quantities
provides a finite upper bound for that coefficient. The claim follows.
There is no need to choose $\epsilon$ uniformly as $m$ varies: the
differentiation is performed for one fixed finite basis at a time.

Let $F(\epsilon)$ be the objective of that basis and $v_\epsilon$ in
degree $m$. By \eqref{eq:center-specialization},
\[
 F(\epsilon)\ge f_{m,t}(y)\ge f_{m,t}(x)=F(0).
\]
Its right derivative at zero is therefore nonnegative. Differentiating
using \eqref{eq:center-discrepancy-variation} and
\eqref{eq:center-order-variation}, and canceling positive common
factors, gives
\begin{equation}\label{eq:center-variation-upper}
 T_m:=\frac{\displaystyle\sum_{(p,\beta)\in\Gamma_m}
                 |\beta|_1e^{-\ell(p)/m}}
                {mQ_m}
       \le\frac dt.
\end{equation}

\subsection{Growth forced by restricted sections}

We now bound the same moment $T_m$ from below. Restrictions to $W$
give many pure tangential initials; multiplication by their sections
forces every nonempty transverse slice to grow with the degree.
For a fixed transverse exponent $p$, define
\[
 G_p(j)=\{\beta:(p,\beta)\in\Gamma_j\},
 \qquad b_p(j)=\bigl(d!\#G_p(j)\bigr)^{1/d},
\]
with $b_p(j)=0$ for an empty slice. The restriction observation above
implies
\[
 \#G_0(h)\ge I_W(h).
\]
Indeed a nonzero restriction has old weight zero and no fiber variables;
elimination within the restriction space gives distinct pure tangential
initials. Lemma~\ref{lem:uniform-restrictions} consequently gives
\begin{equation}\label{eq:center-pure-growth}
 b_0(h)\ge h-C_d
\end{equation}
for all sufficiently large $h$, uniformly in the center and the test.

There is also a useful exact inclusion in bounded degree. Set $h_0=kd$.
Write the selected affine projective coordinates as $s_i/s_0$, with
$s_0(y)\ne0$ and $s_i(y)=0$, for $1\le i\le d$. For
$\beta\in\{0,1\}^d$, the section
\[
 \prod_{i=1}^d s_i^{\beta_i}\,s_0^{d-|\beta|_1}\in R_{h_0}
\]
has initial exponent $(0,\beta)$. The change from the projective
trivialization to our local frame is a unit, whose initial exponent is
zero. Hence
\begin{equation}\label{eq:center-cube}
 \{0,1\}^d\subset G_0(h_0).
\end{equation}

The cube inclusion turns the continuous Brunn--Minkowski inequality
into the discrete estimate needed here. If $A,C$ are nonempty finite
subsets of $\mathbb Z^d$, then
\[
 (A+[0,1]^d)+(C+[0,1]^d)
 \subset A+C+\{0,1\}^d+[0,1]^d.
\]
The volume of a finite lattice set thickened by the unit cube is its
cardinality, since different cubes have disjoint interiors. The
Brunn--Minkowski inequality \cite[Theorem~4.1]{gardner2002}, applied to
these finite unions of cubes, gives
\begin{equation}\label{eq:discrete-bm}
 \#(A+C+\{0,1\}^d)^{1/d}
 \ge (\#A)^{1/d}+(\#C)^{1/d}.
\end{equation}

Apply this with $A=G_p(j)$ and $C=G_0(h-h_0)$. Equations
\eqref{eq:center-semigroup} and \eqref{eq:center-cube} show that
their sum with the cube lies in $G_p(j+h)$. Thus, whenever $G_p(j)$ is
nonempty and $h$ is sufficiently large,
\begin{equation}\label{eq:center-slice-growth}
 b_p(j+h)\ge b_p(j)+h-C'_d,
 \qquad C'_d=C_d+h_0.
\end{equation}
The constants and degree threshold remain independent of $j$, $p$, the
center, and its chosen coordinates.

To sum these slice-growth estimates with the exponential weights, put
\[
 B_j=\frac1{d!}\sum_p e^{-\ell(p)/m}b_p(j)^{d+1},
\]
using the same fixed test and reference degree $m$ as in
\eqref{eq:center-q-b}.
All these sums are finite. For sufficiently large $h$, we have
$h-C'_d\ge0$. Convexity of the $(d+1)$st power and
\eqref{eq:center-slice-growth} yield
\begin{equation}\label{eq:center-b-growth}
 B_{j+h}\ge B_j+(d+1)(h-C'_d)Q_j.
\end{equation}
Only slices present at degree $j$ are needed in deriving this estimate;
new slices make a nonnegative contribution. The coefficient $Q_j$
appears because $b_p(j)^d/d!=\#G_p(j)$.

The power $d+1$ in $B_j$ relates slice size to tangential moment. For any
finite $A\subset\mathbb Z_{\ge0}^d$, let $b=(d!\#A)^{1/d}$. Among
measurable subsets of the positive orthant of volume $\#A$, the simplex
\[
 \{u_i\ge0: u_1+\cdots+u_d\le b\}
\]
minimizes the integral of $u_1+\cdots+u_d$. This follows directly by
comparing the integrand inside and outside this sublevel set. Its
integral is $d b^{d+1}/((d+1)d!)$. On the other hand, the integral over
$A+[0,1]^d$ is $\sum_{\beta\in A}(|\beta|_1+d/2)$. Enlarging $d/2$ to
$d$, applying the resulting inequality slice by slice, and inserting
the weights gives
\begin{equation}\label{eq:center-moment}
 \sum_{(p,\beta)\in\Gamma_m}
 e^{-\ell(p)/m}(|\beta|_1+d)
 \ge \frac{d}{d+1}B_m.
\end{equation}

\subsection{The limiting contradiction}

For each fixed center dimension $d>0$, the preceding bounds hold for
every minimizing pair in the indicated ranges, with slice-growth
constants independent of its center and degree.
We now complete the proof of Proposition~\ref{prop:point-center}. By
Propositions~\ref{prop:ordinary-gap} and~\ref{prop:minimizers}, there is
a common interval of values of $t$ immediately below $n+1$ on which
the minima, for sufficiently large degrees, are bounded away from zero
and one, and their discrepancies lie in a fixed compact interval
$[a_*,a^*]\subset(0,\infty)$. In particular,
\[
 I_*:=\int_0^1 s^n
       \exp\!\left(\frac{(1-s)a_*}{n+1}\right)\,ds
       >\frac1{n+1}.
\]
Choose once and for all $t<n+1$ in that interval sufficiently close to
$n+1$ that
\begin{equation}\label{eq:center-choose-t}
 I_*>\frac1t.
\end{equation}
Keep this choice of $t$ in $f_j=f_{j,t}(x)$. Put
$c=\liminf_{j\to\infty}f_j>0$, and choose
an increasing sequence of degrees $m$ for which $f_m\to c$.

Suppose, seeking a contradiction, that infinitely many of these degrees
admit a minimizing monomial pair with positive-dimensional center.
Choose one such pair in each degree and pass to a subsequence on which
the center dimension is a fixed $d>0$. The centers, valuations, and
models are allowed to vary along this subsequence.

Apply the preceding construction to each chosen pair. Its test and
coordinates may change with $m$, while all auxiliary degrees $j$ for
that $m$ use that one test and coordinate system.

Fix an integer $l\ge2$. For $1\le i\le l$, put
$j_i=\lfloor mi/l\rfloor$, so $j_l=m$. Table~\ref{tab:constants}
summarizes the dependencies and order of choices established above.

\begin{table}[tbp]
\centering
\small
\begin{tabular}{@{}p{.25\textwidth}p{.68\textwidth}@{}}
\toprule
Data & Dependence and order of choice \\
\midrule
$k$, $C_d$, $H_d$, $C'_d$ & Depend only on $(X,L)$ and $d$ after
  fixing a very ample multiple $kL$; independent of the center,
  its degree and singularities, and the test. \\[3pt]
$t_0$, $a_*$, $a^*$ & Fixed by Lemma~\ref{lem:min-coercivity} at the
  hypothetical base point $x$; uniform in all sufficiently large degrees
  and $t\in[t_0,n+1]$. \\[3pt]
$t$, then $m$ & First choose one $t<n+1$ using $a_*$; then take a
  subsequence on which $f_{m,t}(x)$ approaches its positive lower limit.
  Centers and models may vary with $m$. \\[3pt]
$\epsilon$; then $l$ & Differentiate at $\epsilon=0$ for each fixed
  finite basis. In the partition estimate, let $m\to\infty$ at fixed
  $l$, and only afterwards let $l\to\infty$. \\
\bottomrule
\end{tabular}
\caption{Uniform quantities and the order of limits in the center argument.}
\label{tab:constants}
\end{table}

As $m$ tends to infinity on
our chosen subsequence, all increments $j_{i+1}-j_i$ tend to infinity.
Thus \eqref{eq:center-b-growth} applies successively. Starting with
$B_{j_1}\ge0$ and then using \eqref{eq:center-moment} yields
\begin{equation}\label{eq:center-partition-bound}
 T_m+\frac dm
 \ge d\sum_{i=1}^{l-1}
       \frac{j_{i+1}-j_i-C'_d}{m}\,
       \frac{Q_{j_i}}{Q_m}.
\end{equation}
The uniformity of $C'_d$ is indispensable here, since the center and
the model may depend on $m$.

For the auxiliary degrees, the test $(j/m)v$ is admissible at $x$,
because its center is still $W$. Its discrepancy is $j\A(v)/m$, and
its section orders divided by $j$ are $v(s)/m$. Applied to the basis
with initials $\Gamma_j$, the definition of $f_j$ therefore gives
\begin{equation}\label{eq:center-q-compare}
 Q_j\ge N_j f_j e^{-j\A(v)/(mt)}
 \quad\text{for all sufficiently large }j.
\end{equation}
This comparison occurs at $x$ and does not require the chosen point
$y$ to satisfy specialization in degree $j$. Only the degree-$m$
variation used \eqref{eq:center-specialization}.

Combining \eqref{eq:center-q-compare} with
\eqref{eq:center-q-minimum}, we have
\begin{equation}\label{eq:center-ratio}
 \frac{Q_{j_i}}{Q_m}
 \ge\frac{N_{j_i}}{N_m}\frac{f_{j_i}}{f_m}
      \exp\!\left(\left(1-\frac{j_i}{m}\right)
                         \frac{\A(v)}{t}\right).
\end{equation}
For each fixed $i\in\{1,\ldots,l-1\}$, the Hilbert asymptotic gives
\[
 \frac{N_{j_i}}{N_m}\longrightarrow(i/l)^n.
\]
Also $j_i\to\infty$, $f_m\to c$, and $c=\liminf_j f_j>0$, so
\[
 \liminf_{m\to\infty}\frac{f_{j_i}}{f_m}\ge1.
\]
This does not require the degrees $j_i$ themselves to belong to the
subsequence realizing the liminf. Finally, $\A(v)\ge a_*$ and
$t<n+1$ imply
\[
 \exp\!\left(\left(1-\frac{j_i}{m}\right)
                         \frac{\A(v)}{t}\right)
 \ge
 \exp\!\left(\left(1-\frac{j_i}{m}\right)
                         \frac{a_*}{n+1}\right).
\]
All summands in \eqref{eq:center-partition-bound} are nonnegative for
large $m$, and their number is fixed. Taking the liminf therefore gives
\begin{equation}\label{eq:center-riemann-lower}
 \liminf_{m\to\infty}T_m
 \ge \frac dl\sum_{i=1}^{l-1}(i/l)^n
               \exp\!\left(\frac{(1-i/l)a_*}{n+1}\right).
\end{equation}
Only now let $l$ tend to infinity. The right side tends to $dI_*$,
which is strictly greater than $d/t$ by
\eqref{eq:center-choose-t}. This contradicts
\eqref{eq:center-variation-upper}.

There can consequently be only finitely many degrees on the selected
subsequence admitting a positive-dimensional minimizing monomial
center. A nonzero test has a nonempty closed center, and every center
under consideration contains $x$. In all remaining degrees that center
must therefore be exactly $\{x\}$, proving
Proposition~\ref{prop:point-center}.
